\section{Falsifiable route and intermediate claims}
\label{sec:falsifiability}

To avoid a non-testable narrative, we isolate explicit falsifiable intermediate claims. Each item is refuted by numerical experiments or by data confrontation.

\begin{enumerate}
    \item \textbf{Robustness of impedance additivity.} In the three-stratum $\alpha_{\mathrm{em}}$ model, replacing the serial sum by standard alternatives produces large deviations at fixed channel costs (Remark~\ref{rem:alpha_aggregation_isolation}). The same isolation holds for the logarithmic Newton coupling $I_G(p)$ (Remark~\ref{rem:IG_aggregation_isolation}). This isolates the serial rule as the low-complexity composition compatible with the numeric benchmarks.

    \item \textbf{Scale choice and RG/discreteness error budget.} The discrepancy $\alpha_{\mathrm{CODATA}}^{-1}-\alpha_{\mathrm{geo}}^{-1}\approx -3.05\times 10^{-4}$ is mapped to a finite-resolution interface correction in Section~\ref{sec:alpha_anchor}: under one-loop electron-only running the equivalent scale shift is $\mu_{\mathrm{eff}}/\mu_0\approx 1.0029$, corresponding to a sub-step depth correction $\delta r\approx 5.96\times 10^{-3}$ under the Fibonacci map.

    \item \textbf{Resolution-map calibration.} With the golden-branch choice and the Fibonacci map $\mu(r)=\mu_0\varphi^r$ (Axiom~\ref{ax:fib_resolution_map}), the PDG reference scales $\{m_\mu,m_\tau,m_W,m_Z\}$ map to depths within $\delta_0=0.134$ of integers under the rigid reference choice $\mu_0=m_e$ (Table~\ref{tab:resolution_map_calibration}). Under a uniform null for the fractional parts, the corresponding probability is $(2\delta_0)^4\approx 5.2\times 10^{-3}$ (Proposition~\ref{prop:resolution_map_uniform_null}).

    \item \textbf{Electroweak matching rigidity.} The minimal volume-quantized prediction in Theorem~\ref{thm:weinberg_angle} gives $\sin^2\theta_W(\mu_Z)=3/13$ and $\alpha^{-1}(\mu_Z)=13\pi^2$ with the explicit deviations recorded in Section~\ref{sec:running_couplings}. Proposition~\ref{prop:weinberg_rational_rigidity} isolates $3/13$ as the best reduced rational at denominator $\le 50$ for the PDG benchmark.

    \item \textbf{Dispersion signatures.} Discrete scan induces energy-dependent group-velocity corrections (Appendix~\ref{app:c_dispersion}). One must distinguish group velocity from signal-front velocity and preserve causality via a local microscopic model.

    \item \textbf{Mass-ratio rigidity signals.} The proton--electron ratio is fixed by a finite phase-volume assignment (Theorem~\ref{thm:mu_geo}) together with a primitive factorization rigidity check (Table~\ref{tab:mu_phase_factorizations}) and a low-complexity integer rigidity check (Proposition~\ref{prop:mu_integer_rigidity}).

    \item \textbf{CP-odd rigidity signals.} The CKM Jarlskog invariant exhibits a stable low-complexity $\pi$-power signature: Proposition~\ref{prop:jarlskog_rigidity} gives the unique best fit at the PDG central value, and Proposition~\ref{prop:jarlskog_minimax_rigidity} gives the unique minimax solution over the PDG interval in the same bounded search domain.

    \item \textbf{Gravitational coupling in log form.} The logarithmic Newton coupling $I_G(p)=\log(\alpha_G(p)^{-1})$ exhibits a stable low-complexity $\pi$-polynomial signature (Proposition~\ref{prop:alpha_G_rigidity}) and the serial aggregation rule is isolated at fixed three-stratum costs (Remark~\ref{rem:IG_aggregation_isolation}). The phase-space assignment in Axiom~\ref{ax:gravity_log_impedance} is tested directly against CODATA via Theorem~\ref{thm:alpha_G_proton}.
\end{enumerate}
